% Part: computability
% Chapter: recursive-functions
% Section: partial-functions

\documentclass[../../../include/open-logic-section]{subfiles}

\begin{document}

\olfileid{cmp}{rec}{par}
\olsection{పాక్షిక పునరావృత్త ప్రమేయాలు}


పునరావృత్త ప్రమేయాల నిర్వచనానికి ప్రేరణగా, గణనీయమై ఆదిమ పునరావృత్తం
కాని ప్రమేయాలు ఉన్నాయని మన నిరూపణ నిజానికి ఇంకా ఎక్కువ స్థాపిస్తుందని
గమనించండి. వాదన సరళమైనది: $x$, $y$ల ప్రమేయంగా $f_x(y)$ గణనీయం
అయ్యేలా $f_0,f_1,\dots$ ప్రమేయాలను జాబితా చేయవచ్చనే వాస్తవాన్ని
మాత్రమే వాడాం. కాబట్టి \emph{ఇలా జాబితా చేయగల ఏ ప్రమేయాల వర్గానికైనా}
ఆ వాదన వర్తిస్తుంది. ఇది మనలను ఇబ్బందిలో పెడుతుంది: గణనీయ ప్రమేయాలను
స్పష్టంగా వర్ణించాలనుకుంటాం; కానీ గణనీయ ప్రమేయాల సమాహారానికి ఏ
స్పష్టమైన వర్ణనా సమగ్రంగా ఉండదు!

\emph{పాక్షిక} ప్రమేయాలను అనుమతించడమే దీనికి మార్గం. పాక్షిక గణనీయ
ప్రమేయాలను జాబితా చేయడం \emph{సాధ్యమే} అని చూస్తాం.\iftag{TMs}{ నిజానికి,
  ట్యూరింగ్ యంత్రాలను క్రమబద్ధంగా జాబితా చేయవచ్చు కనుక ఇది సాధ్యమని
  మనకు ఇప్పటికే దాదాపు తెలుసు.}{} తరువాత మన వికర్ణ వాదనకు తిరిగి
వచ్చి, పాక్షిక ప్రమేయాలను చేర్చినప్పుడు అది ఎందుకు పనిచేయదో
పరిశీలిస్తాం.

ఇప్పుడు ప్రశ్న ఇది: పాక్షిక పునరావృత్త ప్రమేయాలన్నిటినీ పొందడానికి
ఆదిమ పునరావృత్త ప్రమేయాలకు ఏమి చేర్చాలి? రెండు పనులు చేయాలి:
\begin{enumerate}
\item పాక్షిక ప్రమేయాలను కూడా అనుమతించేలా ఆదిమ పునరావృత్త ప్రమేయాల
  నిర్వచనాన్ని మార్చాలి.
\item కొన్ని కొత్త పాక్షిక ప్రమేయాలు చేరేలా నిర్వచనానికి మరొకదాన్ని
  \emph{చేర్చాలి}.
\end{enumerate}

మొదటిది సులభం. ముందులాగే శూన్యం, ఉత్తరగామి, ప్రక్షేపాలతో మొదలుపెట్టి,
సంయుక్తం, ఆదిమ పునరావృత్తి కింద సంవృతం చేస్తాం. నిర్వచనంలోని కొన్ని
పదాలు నిర్వచితం కాకపోయే అవకాశాన్ని అనుమతించడానికి సంయుక్తం, ఆదిమ
పునరావృత్తి నిర్వచనాలను మార్చడమే తేడా. $f$, $g$ పాక్షిక ప్రమేయాలైతే,
$x$ వద్ద $f$ నిర్వచితమని, అంటే $x$ అనేది $f$ నిర్వచన క్షేత్రంలో ఉందని
సూచించడానికి $f(x)\fdefined$ అని రాస్తాం; దానికి విరుద్ధంగా, $x$
వద్ద $f$ నిర్వచితం కాదని సూచించడానికి $f(x)\fundefined$ అని రాస్తాం.
$f(x)$, $g(x)$ రెండూ అనిర్వచితాలు, లేదా రెండూ నిర్వచితమై సమానాలు
అయినప్పుడు $f(x)\simeq g(x)$ అని రాస్తాం. మరింత సంక్లిష్ట పదాలకూ ఈ
సంకేతనాలను వాడతాం. $h$, $g_0$, \dots,~$g_k$ అన్నీ పాక్షిక
ప్రమేయాలైతే
\[
h(g_0(\vec x),\dots,g_k(\vec x))
\]
అనేది ప్రతి $g_i$, $\vec x$ వద్ద నిర్వచితమై, $h$ అనేది
$g_0(\vec x)$, \dots,~$g_k(\vec x)$ల వద్ద నిర్వచితమైనప్పుడే
నిర్వచితమవుతుందని ఆనవాయితీగా తీసుకుంటాం. ఇలా అర్థం చేసుకున్నప్పుడు,
పాక్షిక ప్రమేయాలకు సంయుక్తం, ఆదిమ పునరావృత్తి నిర్వచనాలు పైనివే;
``$=$''కు బదులు ``$\simeq$''ను పెట్టాలి అంతే.

పాక్షిక ప్రమేయాలను పొందడానికి ఆదిమ పునరావృత్త ప్రమేయాల నిర్వచనానికి
మనం చేర్చేది \emph{అపరిమిత అన్వేషణ క్రియాకారి}. సహజ సంఖ్యలపై
$f(x,\vec z)$ ఏదైనా పాక్షిక ప్రమేయమైతే, $\mu x\;f(x,\vec z)$ను
\begin{quote}
  $f(0,\vec z), f(1,\vec z), \dots, f(x,\vec z)$ అన్నీ
  నిర్వచితమై, $f(x,\vec z)=0$ అయ్యే అతి చిన్న $x$, అటువంటి $x$
  ఉంటే
\end{quote}
అని నిర్వచిస్తాం; లేకపోతే $\mu x\;f(x,\vec z)$ అనిర్వచితమని అర్థం.
ఇది $\mu x\;f(x,\vec z)$ను ఏకైకంగా నిర్వచిస్తుంది.

\begin{explain}
మన నిర్వచనం\iftag{TMs}{ ట్యూరింగ్ యంత్రాలనుగానీ,}{} అల్గారిథంలనుగానీ,
ఏ నిర్దిష్ట గణనా నమూనానుగానీ ప్రస్తావించదని గమనించండి. కానీ సంయుక్తం,
ఆదిమ పునరావృత్తిలాగే, అపరిమిత అన్వేషణ వెనుక ఒక కార్యాత్మక గణనా
అంతర్బోధ ఉంది. పాక్షిక ప్రమేయపు గణనీయతలో, ప్రమేయం అనిర్వచితమైన
ఆర్గుమెంట్లు గణన ఆగని ఇన్‌పుట్‌లకు అనుగుణమవుతాయి. $\mu x\;
f(x,\vec z)$ను గణించే ప్రక్రియ ఇదే: 0 విలువ తిరిగి వచ్చేవరకు
$f(0,\vec z), f(1,\vec z), f(2,\vec z)$లను గణించండి. అయితే
మధ్యలోని ఏ గణనైనా ఆగకపోతే, $\mu x\;f(x,\vec z)$ గణన కూడా ఆగదు.
\end{explain}

$R(x,\vec z)$ ఏదైనా సంబంధమైతే, $\mu x\;R(x,\vec z)$ను
$\mu x\;(1\tsub\Char{R}(x,\vec z))$గా నిర్వచిస్తాం. మరో విధంగా,
$R(x,\vec z)$ నెరవేరే అతి చిన్న $x$ను $\mu x\;R(x,\vec z)$ తిరిగి
ఇస్తుంది. కాబట్టి $f(x,\vec z)$ సర్వనిర్వచిత ప్రమేయమైతే,
$\mu x\;f(x,\vec z)$ అనేది $\mu x\;(f(x,\vec z)=0)$తో సమానం.
కానీ మన అసలు నిర్వచనం మరింత సాధారణమని గమనించండి: అది $f(x,\vec z)$
ప్రతిచోటా నిర్వచితం కాకపోయే అవకాశాన్ని అనుమతిస్తుంది. దీనికి విరుద్ధంగా,
ఒక సంబంధపు లక్షణ ప్రమేయం ఎల్లప్పుడూ సర్వనిర్వచితమే.

\begin{defn}
సహజ సంఖ్యల నుంచి సహజ సంఖ్యలకు వెళ్లే వివిధ స్థానసంఖ్యల పాక్షిక
ప్రమేయాల్లో, శూన్యం, ఉత్తరగామి, ప్రక్షేపాలను కలిగి, సంయుక్తం, ఆదిమ
పునరావృత్తి, అపరిమిత అన్వేషణ కింద సంవృతమైన అతి చిన్న సమితిని
\emph{పాక్షిక పునరావృత్త ప్రమేయాల} సమితి అంటాం.
\end{defn}

పాక్షిక పునరావృత్త ప్రమేయాల్లో కొన్ని ప్రతి ఆర్గుమెంటు వద్ద
నిర్వచితమై, సర్వనిర్వచిత ప్రమేయాలు అవుతాయి.

\begin{defn}
\ollabel{defn:recursive-fn}
సర్వనిర్వచితమైన పాక్షిక పునరావృత్త ప్రమేయాల సమితినే
\emph{పునరావృత్త ప్రమేయాల} సమితి అంటాం.
\end{defn}

పునరావృత్త ప్రమేయం ప్రతిచోటా నిర్వచితమని నొక్కి చెప్పడానికి దాన్ని
కొన్నిసార్లు ``సర్వనిర్వచిత పునరావృత్త ప్రమేయం'' అంటారు.

\end{document}
